% Authorship: Claude-drafted, author-revised
\section{Computational Framework}

This Part explains how Berman's theory is implemented as testable
computational models and sets out the four requirements against which those
models are measured. Already formal enough to implement directly, the theory admits of the
systematic variation that thought experiments cannot supply. Three sections
follow. The first Section describes the methods. At its core are two things: a direct rendering of Simple Additive
Weighting, the \textit{stipulated engine}, whose weights are set by hand,
and four mechanisms that ground the principle weights through four
different practices. In order, the four practices are the uptake of principles in judicial
reasoning, their reliance-based embedding in precedent, the
strategic adaptation of judges, and the exogenous shock of judicial
appointment.

By the one line the theory is defined against, the second Section then
sorts the candidate aggregation functions: whether a rule weighs the
magnitude of competing principles or merely ranks them. The third states the four requirements any faithful implementation must
meet. On that test the weights must be well-defined, tolerant of rough
specification, magnitude-weighing, and productive of more law than a Hartian baseline. On the
evidence reported there, every mechanism meets them.\footnote{All code, data, and analysis scripts are
publicly available. See
\url{https://dgk-law-and-cognition-lab.github.io/BermanModeling/}. The
implementation is in Python with standard scientific-computing libraries.
Case and seed counts vary by experiment rather than following one convention.
Synthetic case draws run from three hundred to fifty thousand, and repetitions
from a single seeded draw to thirty independent seeds. The four requirements
stated below are scored on eight hundred held-out cases across thirty seeds
for each of the four mechanisms. Confidence intervals are 95\% bootstrap
intervals from 2{,}000 resamples, computed from a fixed seed so that every
reported interval is exactly reproducible. Most comparisons in what follows
are reported descriptively, as means across seeds rather than as hypothesis
tests. Two are tested: the comparison across replacement rules in Part~V uses
the Wilcoxon signed-rank test on the seed-paired values, and the same Part
tests unimodality with Hartigan's dip test, whose $p$-values are interpolated
from a simulated null rather than exact.}

\subsection{Methods}

\subsubsection{The Stipulated Engine}

The core implementation encodes Berman's aggregation equation directly. In it
each principle carries a name, a weight $w_i \in [0,1]$, and a category. Nine principles run throughout, and the assignment of each to
a category is the \textit{category map}: text and original meaning are the two textual principles; precedent,
structure, and federalism the three structural ones; and dignity, equality,
liberty, and welfare the four moral ones. They
track the examples in Part II. The three categories are of unequal size, which matters later for any rule
that takes them in order. Such a rule calls them tiers, as the hierarchical
Hart does, and this Article uses the two words for the same three groups. A case maps each principle to an activation
$a_i \in [-1,+1]$; net force is $F = \sum_i w_i a_i$; and a case is
determinate when $|F|$ exceeds the threshold $\tau$. As Part II argued, the threshold is a modeling choice rather than a
discovered constant. The faithful baseline is the bare sign of net force
($\tau = 0$). As Part VIII shows, a positive $\tau$ is a confidence gate, not a quantity
the theory fixes. Unless a passage says otherwise, the results below are computed at $\tau =
0.2$, a gate this Article imposes rather than one the theory supplies. Part~IV
reports the fate of conflicting cases at the bare sign as well, and
Part~VIII sweeps the threshold from $0$ to $0.30$ to separate the results that
move with it from those that do not.

Throughout, I test the model on synthetic cases rather than hand-coded
disputes.\footnote{Validation on real disputes is the natural next step but
requires extensive doctrinal coding (assigning activation values to actual
judicial opinions) that is beyond this Article's scope. Part X sketches what
such an empirical program would involve.} Unless otherwise noted, activations are drawn independently from a uniform
distribution on $[-1,+1]$.\footnote{How the mechanisms generate their own case streams departs from this, and the
departure is set out here rather than left to be found in the code. The uptake
bench draws its stream independently, as stated. Reliance, strategic
adaptation, and appointment generate their weights on a category-correlated
stream instead, each activation mixing a latent factor shared by the
principles of its category with idiosyncratic noise, weighted $0.7$ to the
factor. A third structure, \textit{full correlation}, which later Parts use as
a limiting case, replaces the per-category factors with a single latent factor
shared by all nine principles at the same $0.7$ weighting, so that all nine
tend to move together regardless of kind. The four requirements are then scored for all four mechanisms on the
same independent held-out battery, so the pass-or-fail verdicts of the
requirements Section rest on the conservative baseline throughout. The
generating stream nonetheless shapes the weights themselves. Run on an
independent stream, reliance keeps the direction of its lean toward the moral
principles but loses most of its size, the gap between its mean moral weight
and its mean textual weight falling from about $0.40$ to about $0.04$ on
weights rescaled so the heaviest principle is one. The corresponding
comparison has not been run for appointment.} This is a conservative choice. Independent
activations maximize variance and conflict. By contrast, real cases would likely exhibit correlated activations, which
Part IV shows improve the model's noise tolerance and its raw determinacy.
On independent activations, therefore, those results understate how the
model would fare on realistic ones. As Part IV reports, the margin over the hierarchical Hart is the exception
and runs the other way; the requirement is accordingly scored on independent activations, where the
comparison would seem less flattered by a tiered rule's ability to stop at
the top tier.

Rather than fix one set of weights, I test three distributions spanning
plausible weight structures.\footnote{The three are \textit{uniform} ($w_i$
drawn from $U(0.60, 0.70)$, near-equal importance), \textit{moderate} ($w_i$
from $U(0.50, 0.80)$), and \textit{high-variance} ($w_i$ from $U(0.30,
0.95)$, a system in which some principles matter much more than others).}

To measure how much weight imprecision the model tolerates, I perturb each
weight with Gaussian noise and ask how often the perturbed verdict matches
the original. From that perturbation comes the key quantity, the \textit{critical
uncertainty} $\sigma^{*}$: the noise level at which agreement with the unperturbed verdict
falls below eighty percent, the point at which one case in five changes its
result under perturbation.\footnote{Formally, I add $\epsilon \sim N(0,
\sigma^2)$ to each weight, clipped to $[0,1]$. A case counts as robust at a
given $\sigma$ if it keeps its unperturbed \textit{signed} verdict
(determinate-for, determinate-against, or indeterminate) on at least ninety
percent of forty random redraws of the noise (Monte-Carlo draws); $\sigma^{*}$ is the smallest
$\sigma$ at which the share of robust cases drops below 0.80. Tracking the
signed verdict is stricter than tracking determinacy alone (a flip from
determinate-for to determinate-against counts as instability), and the
search grid stops at $\sigma = 1.0$, so a model still robust there is recorded
as tolerant beyond it. The correlation comparison Part~IV reports is the one
departure from the count of forty, using sixty-four; since its arms share that
count, the contrast within it is unaffected.} Throughout, two further quantities appear as descriptive measures, never as
pass-or-fail thresholds: the \textit{coefficient of variation} of a weight
vector, recording how far its entries lie from their own mean, and the
number of \textit{surviving} principles, those carrying weight above a small
cutoff. The first of these has to be read with care where a mechanism sets
some weights to zero, since a vector holding zeros disperses widely without
being graded among the weights it retains.

\subsubsection{Four Mechanisms of Practice}

The stipulated engine takes the weights as given. But the weights are the
theory's hinge. Berman holds, correctly on this Article's account, that
they are conferred by practice rather than stipulated by a
theorist.\footnote{Berman, \textit{Practices}, \textit{supra} note 7, at
329.} The open question is which practice, and through what dynamic. The next four models answer it differently.

Each grounds the principle weights through a different practice:
uptake, reliance, strategic adaptation, and appointment. In running all
four, the purpose is to show that the formal apparatus does not depend on
which is right. As the requirements Section details, each yields weights that meet them, on every seed for three of the four mechanisms and on all but
a handful of seeds for the fourth. The four do not share the weights themselves, and that is where the
theory's open question lies.

\textit{Uptake.} The first mechanism grounds weight in how often a principle
is taken up. On this model a bench of roughly one hundred judges decides a stream of
cases in nine-judge panels. Each judge carries one of four interpretive methodologies, which sets the
larger part of its propensity to invoke each principle. Two smaller terms
adjust it: a judge leans slightly toward the principles the bench currently
weighs most, and reaches somewhat more readily for a principle the case before
it activates strongly. From that stream emerges each principle's \textit{uptake share}: the
frequency
with which the bench invokes it, expressed as a fraction of the most-invoked principle's own frequency, so that the most-invoked principle is
$1.0$.\footnote{The
uptake share is a spare model of Berman's notion of
``taking-up.'' Berman, \textit{Practices}, \textit{supra} note 7, at 329. It
is not yet weight in his fuller sense, which the reliance mechanism
approaches more closely.}

That share is not yet the weight the model runs on. Between the two sits a
\textit{live/dead cut}: a principle whose uptake share falls below the median
share is set to zero, and the principles at or above the median keep the
shares they have.\footnote{The cut is a different operation from the small cutoff used to
count surviving principles, which is fixed at $0.05$ and only reads a weight
vector rather than altering one. The two are not independent: the cut sets a
weight to zero, and a zeroed weight then falls below the counting cutoff,
which is why the survivor count tracks the cut. The cut is distinct too from the determinacy
threshold $\tau$, which gates verdicts rather than weights.} The cut
implements the threshold of uptake described in Part~II, the point at which
enough officials invoke a norm for it to count as a principle at all. On nine
principles it leaves five standing, the same five on every seed, so the
survivor count reported later is a product of the cut and not of the uptake
shares alone. And it operates while the simulation runs rather than on a
finished vector, so a principle once zeroed is invoked less often in the cases
that follow, making the two-group structure Part~V reports in part a
structure the cut helped produce.

Because one might worry that the result turns on
how the bench changes over time, I run four turnover regimes. For the headline result the bench is static, with no replacement. The others replace by appointment, reseating a random fifth of the bench
with judges drawn afresh from the four methodologies in equal proportion; by
drift, reseating a random fifth with near-copies of judges already
sitting; and by a fitness rule that retires the least decisive judges.
Replacement by appointment is a turnover rule inside this bench, not the
appointment mechanism described below, which confers weight by a different
measure entirely. Under all four the emergent structure holds. Notably, it does not depend on the fitness rule.
Berman has questioned that rule's Darwinian framing, and the design is built
to do without it. On this model, a principle's weight tracks how often a practice
takes it up, among the principles the cut leaves standing.

\textit{Reliance.} The second mechanism is the closest to Berman's own
account: weight conferred by \textit{embeddedness}. A principle gains weight by prevailing in contested, dispositive disputes.
How much the resulting holding is later relied upon amplifies that gain,
which decays over time and is stripped away when a precedent is overruled.\footnote{Decision and
read-out are decoupled. The disputes are resolved by a fixed, exogenous
pressure structure, and the conferred weight is read out separately, so the
dynamic cannot lock itself in. A circular variant, in which conferred weight
feeds back into the disputes it decides, is implemented only to demonstrate
the runaway lock-in it would cause, and is switched off in the reported
runs.} When Berman says that a principle is preserved by the holdings that rely on
it, this is arguably the mechanism he gestures at. From uptake, which counts invocations and reads nothing about how a dispute
came out, it is distinct, and distinct too from the strategic mechanism below,
which optimizes toward a goal. Here a principle's weight
is the accumulated, decaying weight of the precedent that rests on it.

\textit{Strategic.} The third mechanism models judges as adaptive
optimizers.\footnote{See generally Richard S.\ Sutton \& Andrew G.\ Barto,
\textsc{Reinforcement Learning: An Introduction} (2d ed.\ 2018).} Each is a reinforcement-learning agent that adjusts a shared weight vector
to maximize a reward. That reward is consensus agreement routed through the weighted sum and
computed on a slightly noised draw of each judge's force, so that the weights
receive a gradient at all. It rewards agreement, adds a bonus where the agreed
force is decisive, and charges a small cost for each principle invoked. That
cost keeps the optimum from being to invoke everything.\footnote{The
quantity read out is a genuine weight that multiplies the activations inside
each judge's force, not a probability of invoking a principle. Judges also
learn when to invoke, in the same update, but that policy is not what the
weights are read from. (An earlier version of this project read out the
learned invocation frequencies instead, which merely reproduced uptake; that
is corrected here.) Learning uses a standard reinforcement-learning update, the REINFORCE
policy gradient: each agent nudges its weights toward whatever improved its
reward, and the nudges shrink toward zero over training so the weights
settle to a stable read-out.} Of the four this is the least Bermanian. Berman nowhere describes judges
optimizing a goal, and the mechanism should be read as a robustness lens
rather than as an interpretation of the theory. If even strategic learning
produces weights that meet the requirements, those requirements are
unlikely to be an artifact of any one picture of practice.

\textit{Appointment.} The fourth mechanism builds on the dynamic Berman himself proposed: the
near-random, exogenous shock of judicial appointment. The decision pressure that resolves disputes comes from the sitting bench.
At intervals a random cohort of judges is replaced by newcomers, each carrying
an idiosyncratic variation on the common legal culture's pressure structure,
meaning the fixed pattern of decision pressure that settles disputes from
outside the weight dynamic, the same exogenous structure the reliance
mechanism takes as fixed.
They are not selected for success, simply reseated,
which models the social fact that benches turn over.\footnote{There is no fitness
anywhere in this mechanism: the cohort is replaced at random, reflecting
social acceptance rather than selection.}

Weight is read from \textit{holding-support}: the breadth of still-live
holdings resting on each principle, lightly decayed. A holding stays live
until it is overruled, slowly and case by case, or until one of its
principles is disavowed, quickly and wholesale. On this construction a
principle's vitality is its surviving holding-support, which does not
require that the principle be currently invoked. That is the embedding-via-holdings point Berman has emphasized in correspondence. This
mechanism joins his proposed appointment dynamic to that embedding account.

An earlier version of this project also tested a hybrid mechanism that
combined learning with periodic replacement. It is omitted here. Originally it was introduced to find a mechanism that could jointly satisfy
properties an earlier framing scored pass-or-fail. As the requirements Section explains, two of those properties should not
have been scored at all, which
removes the problem the hybrid was built to solve.

Through four different practices, then, these mechanisms ground the
weights. Yet as the requirements Section reports, each yields weights that are
well-defined, noise-tolerant, magnitude-weighing, and more determinate than the
hierarchical Hart, and each does so on every seed but for one mechanism's
noise-tolerance. The formal apparatus is robust to how practice is modeled.
The four do not agree on the weights themselves, and Part VI takes
up the consequences. The aggregation works under every mechanism. At issue,
then, is not the aggregation but the dynamic of changing grounds that
confers the weights it works on.

\subsection{Aggregation Functions}

Parts IV through VI take Simple Additive Weighting as the aggregation
function, the form Berman writes down. But SAW is one member of a family,
and the family is defined by what the theory is set against. Berman places principled positivism on the \textit{weighing} side of a
line. Principles contribute more or less force according to how strongly
they apply. They do not rank in a fixed priority under which a higher-ranked principle would defeat any lower-ranked one regardless of magnitude.\footnote{Berman,
\textit{Practices}, \textit{supra} note 7, at 311--13, 333.}

Sorted by that line into three classes, ten aggregation rules go to Part
VII: the cardinal family that reads magnitudes, the ordinal rules that read
only signs and ranks, and the lexical rules on the far side. The aim is not a
contest for the best rule. It is to confirm that the defining property holds
across the family Berman endorses, and to measure what the choice between
weighing magnitudes and merely counting them is worth.

Of the ten, four are the genuine candidates, all of them cardinal rules that
read magnitudes. \textit{Simple Additive Weighting}, the reference form,
is $F = \sum_i w_i a_i$, determinate when $|F| > \tau$. \textit{Nonlinear combination} keeps the additive form but lets a
principle's contribution grow faster than linearly in how strongly it
applies, $F = \sum_i w_i a_i e^{\beta |a_i|}$ with $\beta =
4$.\footnote{This generalizes SAW while respecting Berman's commitment that
principles contribute weighted force. It is the sharpest departure from linearity in the cardinal family and the
highest scorer on the magnitude probe described below, which is what earns it
a place: it tests whether the defining
property depends on the contribution being linear in activation. It adds no
interaction between principles, so it is not a test of whether principles
reinforce or dampen one another. No rule in the battery tests that.} Holding the form fixed, two further variants vary its surroundings: a
\textit{bare-sign} rule, deciding on the sign of net force, and a
\textit{normalized} rule, rescaling the weights to sum to one before
adding. The normalized rule is a control rather than a fifth function, since a
positive rescaling cannot change the sign of a sum and it therefore returns
the bare-sign rule's verdict on every case; it is carried to forestall the
objection that the results turn on the scale of the weights. Of the four
cardinal rules only Simple Additive Weighting carries a determinacy gate. The
other three decide on the bare sign and abstain on no case, which is why
Part~VII compares them on verdicts rather than on how much they settle. Part VII shows that the in-theory results are not specific to
Simple Additive Weighting: the defining property belongs to every rule in the
family.

Two functions are the mode the theory is defined against. \textit{Lexicographic
ordering} ranks the principles by weight and lets the highest-ranked active
principle decide, lower-ranked principles having no effect whenever a
higher-ranked one is active.\footnote{Lexicographic ordering is associated
with Joseph Raz's view of exclusionary reasons. See Joseph Raz,
\textsc{Practical Reason and Norms} 35--48 (2d ed.\ 1999).}
\textit{Cascade} ranks the principles by category first, text before
structural principles such as federalism and those before substantive ones
such as dignity, and by weight within a category. The highest-ranked principle
that applies at all decides the case, and nothing below it is consulted. Since
activations are drawn from a continuous distribution, a textual principle
applies in effectively every case, so cascade in practice decides on text
alone.\footnote{On twenty thousand cases at the moderate weight vector, the
deciding principle was the plain text in every one. That is a property of the
case generator rather than of the rule, and it is what makes cascade the more
extreme of the two lexical foils.}\footnote{Cascade is the functional form
implicit in categorical textualist orderings. See Antonin Scalia \& Bryan
A.\ Garner, \textsc{Reading Law: The Interpretation of Legal Texts} 56--77
(2012).} These are not candidate readings of principled positivism
entertained and then rejected on performance. They are the structural
contrast class, and the magnitude-weighing requirement measures how far the
weighted sum departs from them.

Four rules isolate the cardinal question. Each keeps the sign, and at most the
within-case rank, of the activations while discarding their magnitudes: a bare
\textit{vote-count}, a \textit{weighted vote} in which a principle's vote
counts in proportion to its standing, a \textit{Borda count} awarding each
principle points equal to its within-case rank in force, scaled by the rank of
its weight rather than by the weight itself, and a \textit{median} rule that
takes the middle vote.

Only one of the four is a clean one-variable contrast with the weighted sum.
The weighted vote keeps the cardinal weights and discards only the magnitudes,
so it isolates exactly the question the magnitude-weighing requirement asks. The
bare vote-count and the median rule drop the weights as well, and the Borda
count replaces them with their ranks, so those three differ from the weighted
sum in two respects rather than one. The median rule is in any event a control
rather than a fourth tally: on a panel of odd size, and the panel here is nine
principles, the median sign is the majority sign, so it returns the
vote-count's verdict on every case. Counting the two controls, the ten labels
name eight distinct verdict rules.

The ordinal rules are not a foil but a rival. Nothing in a legal system
that fixed content by a tally among principles would be incoherent. And
Berman himself allows the principles to carry equal weight, the case
where a weighted sum comes closest to a count.

That is arguably
what makes the cardinal-ordinal line the one worth measuring. Part VII sorts the ten rules by the magnitude
probe, reports that the defining property belongs to the cardinal family alone,
and measures what the choice across that line is worth in verdicts. Across the classes it does not read determinacy rates or Hart margins as
measures of fidelity, because the ungated rules turn every case into a verdict
and no such comparison would be like for like.\footnote{The Borda count is the
one partial exception, abstaining on the small share of cases where its rank
sums tie exactly, which is $0.36\%$ of the battery. Part~VII reports the
inflated cross-class figures and sets them aside on this ground.}

\subsection{Four Requirements}

To compare the mechanisms and functions, I hold each to four
requirements, properties any faithful implementation of principled
positivism must have. Each is drawn from something Berman claims or commits
to, and each is scored seed by seed across thirty random seeds for each of the
four mechanisms. A mechanism counts as meeting a requirement when it clears it
on all thirty; where it clears fewer, the seed count is reported and the
shortfall noted.\footnote{The implemented rule is a seed share rather than
unanimity, requiring $0.90$ of seeds on the first, third, and fourth
requirements and $0.80$ on the second. On the runs reported here the slack
never decides an outcome: three mechanisms clear all four requirements on
every seed, and the fourth falls short only on noise tolerance, at twenty-six
of thirty, which the requirements Section reports rather than absorbs.} The
aggregation functions of Section III.B are held to one of the four rather than
all of them, since only the magnitude probe described below can be run like
for like
across rules that differ in whether they abstain at all; Part~VII explains
why. The four are independent: meeting one neither entails nor precludes
meeting another. Table~\ref{tab:requirements} summarizes them;
the paragraphs that follow give each its basis and its operational test.

\begin{table}[ht]
\centering
\small
\caption{The Four Requirements}
\label{tab:requirements}
\begin{tabular}{@{}p{3.4cm} p{5.2cm} p{4.3cm}@{}}
\toprule
\textbf{Requirement} & \textbf{Operational test} & \textbf{Basis in Berman} \\
\midrule
Well-defined & At least two principles carry weight above a small cutoff.
& The theory's plurality of weighted principles. \\
\addlinespace
Noise-tolerant & Critical uncertainty $\sigma^{*} \geq 0.07$: the weights
survive $\pm 0.07$ of perturbation before one case in five flips its
verdict. & ``Rough weights suffice,'' his answer to Perry. \\
\addlinespace
Magnitude-weighing & Under a magnitude-only perturbation, the weighted sum
changes
its verdict more than five times as often as a lexicographic rule on the
same weights. & Principles weigh; they do not rank. \\
\addlinespace
More law than Hart & A positive determinacy advantage over the
hierarchical Hart baseline. & His answer to the ``too little law''
worry. \\
\bottomrule
\end{tabular}
\end{table}

The first requirement is that the weights be \textit{well-defined}. Unless at least two principles carry genuine weight, the aggregation
formula is vacuous. A model whose weights collapse to nothing, or onto a
single dominant principle, is not testing Berman's theory of a plurality of
weighted principles.\footnote{Berman, \textit{Practices}, \textit{supra} note 7, at
322--24.} Operationally, at least two principles must carry weight above a
small cutoff.

This is a minimum, not an achievement. Its virtue is that it cannot be gamed,
and clearing it earns no credit toward the other three requirements. Nor
should it be confused with a demand that the weights be spread out or
graded, a separate and non-required property discussed below.

The second requirement is that the weights be \textit{noise-tolerant}: small
errors in how heavily each principle counts must not flip the verdict in too
many cases. This is Berman's answer to Stephen Perry, who objected that
aggregation demands a precision that messy practice cannot supply; Berman
replies that rough weights suffice.\footnote{Berman, \textit{Practices},
\textit{supra} note 7, at 328--30; see Perry, \textit{supra} note 9, at 794.} Made quantitative, the requirement is that the weights tolerate at
least $\pm 0.07$ of Gaussian perturbation before more than one case in five
changes its verdict, that the critical uncertainty $\sigma^{*}$ be at least
$0.07$.

The floor is not arbitrary, and its provenance is worth stating exactly. It is
set from the stipulated engine's own tolerance across a grid of
forty-five configurations, three weight distributions by three principle
counts by five thresholds, each cell a thirty-seed mean. Those cell means run
from $0.071$ to $0.134$ and average $0.090$, and the lowest of them falls at
nine principles, moderate weights, and $\tau = 0.2$, the configuration
this Article uses throughout. The floor sits just below that
lowest cell mean.

The floor does not bound the engine from below on every run, and
the difference matters because the requirement is scored seed by seed while
the grid reports means. On larger samples of the same engine at the same
configuration, the per-seed critical uncertainty falls below $0.07$ on
eighteen of thirty seeds in one run and twenty of thirty in another, and the
lowest cell mean's own confidence interval runs from $0.068$ to $0.073$, so
the floor sits inside it.\footnote{The two runs are the eight-hundred-case
correlation baseline and the thousand-case magnitude-probe run, at mean
$\sigma^{*}$ of $0.069$ and $0.068$ respectively. Neither is a defect in the
engine; both are the same quantity measured on more cases, where seed-to-seed
variation is visible rather than averaged away.} So the requirement asks an
emergent mechanism to be about as noise-tolerant as the stipulated engine is
on average at its least tolerant setting, a weaker benchmark than
asking it to beat the engine's worst run, and it should not be read as the
stronger one. It remains internal to the model rather than externally
stipulated, and a lower floor would only make it easier to clear.

The third requirement has no counterpart in earlier treatments of the
theory, and it is perhaps the one that most directly captures what makes
principled positivism the theory it is. The model must be \textit{magnitude-weighing}: it must respond to how strongly
each principle applies, not merely to which principle ranks highest. Berman is explicit that his principles weigh rather
than rank, contributing more or less force rather than standing in a fixed
order of priority.\footnote{Berman, \textit{Practices}, \textit{supra} note
7, at 311--13, 333.}

The operational test is the \textit{magnitude probe}, used throughout the
Article under that name. Holding every sign and every within-case rank fixed,
it perturbs only the magnitudes of the activations and asks how often the
verdict changes. A rule that
merely ranks is invariant to such a perturbation; the weighted sum is not.
On that probe the requirement is that the weighted sum change its verdict
more than five times as often as a lexicographic rule does on the same
weights.

Two
things here are worth separating. That the lexical rule scores near zero is
partly definitional: the probe holds the ranks fixed, and a rule that
consults only ranks is blind to it. What is not definitional is that the
weighted sum responds substantially. Under a magnitude-only perturbation the
ungated sum changes its verdict at roughly ten times the lexical rule's
residual rate, and the gated sum this Article runs at roughly twenty times,
the gate adding the cases a magnitude change pushes across the determinacy
threshold rather than across zero.\footnote{Scores of $0.0256$ ungated and
$0.0502$ gated against the lexical rule's $0.0026$, thirty-seed means. The
requirement is scored on the gated rule, so the ratio it actually applies is
the larger one; the ungated ratio is reported alongside it because it isolates
changes of direction alone.} It is that responsiveness the requirement
credits. Because the five-fold bar sits well below either gap, the bar is a
conservative one, perhaps markedly so.

The fourth requirement is that the model deliver \textit{more law than
Hart}: that it decide strictly more cases determinately than a Hartian
baseline that lets the law run out wherever the principles it consults fail to
converge. This is Berman's answer to the Dworkinian worry that positivism yields too
little law. Where a Hartian baseline finds a gap, summing
weighted forces finds law.

Against the hierarchical Hart the requirement is adjudicated: that is the
charitable baseline, sorting the principles into tiers by category and
deciding a case at the highest tier holding any active principle, and only
where enough of that tier's active principles point the same way. Where they
do not, the law runs out, because the tiers below are never reached. Nothing
in it ranks principles by weight. That separates it from the
single-principle rule the third requirement uses as its foil. On that
comparison the model's
advantage is about $+0.27$, and positive on every seed. Its size is a finding, not a fixed bar, and it would presumably move with
the case distribution.

Against a consensus
Hart, which decides only where enough of the principles in play point the
same way, the advantage is reported as a curve rather than scored. It runs from
slightly negative ($-0.02$) at a fifty-percent consensus to $+0.27$, $+0.59$, and $+0.79$ as
the consensus bar rises to sixty, seventy, and eighty percent. Deliberately,
the single negative cell is included. It arises at the competitor's loosest
bar, where the rival decides more cases than at any other setting on this
battery, and not at its most demanding. It is not a failure of the
requirement, which is scored against the hierarchical Hart.\footnote{The Hart baseline's tiers come from
the exogenous category map, textual on top, structural below, moral below
that, never from any emergent weight, so the comparison does not import the
model's own output into its competitor.}

Two properties that earlier treatments folded in as requirements are
measured here but not scored. One is weight \textit{gradation}, whether the weights are spread out rather
than clustered. On it the mechanisms differ sharply, though not in the way a
single figure suggests. The uptake weights have a coefficient of variation
near $0.90$ and the reliance weights near $0.24$. The larger figure,
though, records the live/dead cut rather than gradation: it measures the
distance between five live principles and four zeros. Among the five that
survive, the uptake weights run from $0.79$ to $1.00$ and their coefficient of
variation is $0.09$, making them the most tightly clustered of the four
mechanisms rather than the most spread. Reliance applies no cut, so its $0.24$
is dispersion among nine live weights and the two figures do not measure the
same thing. Both are reported descriptively.
The other is \textit{correlation benefit}, whether performance improves under
the correlated activations of realistic cases, as it does.

Because Berman's theory demands neither, neither is a requirement: he
permits principles to carry equal weight,\footnote{Berman, \textit{Practices},
\textit{supra} note 7, at 323--24 (the account does not require that the
principles have varied weights).} and he nowhere claims his model needs
correlation to work. The demotion changes the attribution to Berman, not the
roster of surviving models: no mechanism's verdict on the four requirements
turns on it.

Held to these four requirements across thirty seeds, every mechanism
passes, one of them only barely. Part VI reports the margins and what the
near-miss shows. The requirements are therefore co-satisfiable, and they
are so because the verdicts rest on the four alone. Gradation and
correlation benefit, measured but not scored, carry no pass-or-fail weight.
Differences among mechanisms on those two descriptive measures therefore
leave every survival verdict untouched.

These four requirements are neither arbitrary nor exhaustive. Others could
be specified (stability under changes to the case base, robustness to
adversarial manipulation, agreement with human intuitions about particular
cases), and the analysis would probably be richer for testing them. The four here seem the ones most directly tied to Berman's own claims and to
the structural needs of aggregation.

Two limits bear emphasis. The cases
are synthetic. Whether real adjudication grounds its weights through any of
these mechanisms is an empirical question this Article frames but does not
resolve, and one that would likely require coding actual judicial opinions
to answer.
And the determinacy threshold is a modeling choice; the analysis reports
the model's sensitivity to $\tau$, which
accounts for most of the variation in $\sigma^{*}$, rather than treating any value of $\tau$ as canonical. Results reported at
one $\tau$ should not be read as holding at another.

The framework is now in place. Part IV turns to the stipulated engine and
establishes that the weighted sum works well when given faithful weights, and that several of its properties are not contingent findings but
theorems. The Parts that follow then ask where those weights come from, likely the
harder question.
